\section{Analytical Components of Agentic Visual Generation}
\label{sec:components}

The preceding section defined agentic visual generation as state-dependent control over a visual creation trajectory. This section examines the components through which that control is organized during visual creation. A typical trajectory contains a task requirement, a selected operation, an artifact or execution result, an inspection or check, and a subsequent decision, among other task-relevant events. The decision may keep the current plan, change an input or tool, repair part of the artifact, replan, stop, or take another permitted control action. One controller may perform these steps, or different components may handle planning, execution, inspection, and related functions.

For cross-paper comparison, we use six descriptive components. Goal Understanding, Specification, and Planning cover the task goal, constraints, dependencies, and planning records that guide later operations. Memory covers information and records retained for access at later decision points. Tool covers callable models, editors, code environments, applications, simulators, and related interfaces, together with their input conditions and returned execution information. Perception covers signal acquisition, observation, interpretation, diagnosis, verification, and acceptance evidence. Action covers artifact transformation, information acquisition, execution, environment interaction, communication, coordination, trajectory control, recovery, stopping, and related operations. Cross-Task Self-Improvement covers how information from a completed trajectory is carried into a later task: episodic reuse retrieves earlier attempts or repairs, while cross-task self-improvement requires evidence that the update changes later behavior. These headings provide a common record for the functions and records examined in the papers.


\subsection{Six-Component Analytical Framework}
\label{sec:analytical-framework}


\Needspace{7\baselineskip}

\subsubsection{Goal Understanding, Specification, and Planning}

This component interprets task information, expresses the intended visual outcome and its constraints, and organizes the remaining work into a plan. The input may combine a request, references, interaction context, domain requirements, and other task-relevant information with different levels of detail. The component turns this material into commitments that later operations can use, while keeping unresolved information attached to the part of the task that it may affect. Its direct outputs include a current operational specification, unresolved items that affect that specification, and planned subgoals with their dependencies and order. When user information or task evidence changes a commitment or dependency, the affected specification and the remaining plan can be updated.

\paragraph{Goal understanding.}
Goal understanding identifies the intended outcome, the information that bears on it, and the parts of the request that remain unresolved before planning. It uses information including user instructions, references, interaction context, domain requirements, and other task-relevant material to form a clarified task target. The interpretation separates information supplied by the user from information inferred from context and records which unresolved items may change later work. Unresolved information can lead to a clarification request, a scoped assumption, a deferred item, or another task-appropriate decision. A clarification request can obtain missing information, a scoped assumption can support provisional planning, and a deferred item can remain attached to the commitments that depend on it. This association gives later planning a place to apply the new information when the item is clarified or an assumption is revised. Goal understanding therefore supplies the current task interpretation from which the operational specification is constructed.

From Idea to CAD uses interactive requirements elicitation to clarify design intent from sketches and text before construction \cite{ocker2025ideatocad}. Search Beyond What Can Be Taught studies generator-specific knowledge gaps and trains a reasoner to decide when external search is useful; SearchGen-Bench evaluates search-augmented generation, while the co-training experiments study coordination between the search reasoner and the generator \cite{wang2026searchgen}. Clarify Before Executing detects incomplete or ambiguous 3D requests and asks targeted questions before tool orchestration, progressively turning the clarified intent into an executable workflow \cite{avg260716352}.

\paragraph{Operational specification.}
An operational specification is the task-level representation that turns the goal, references, constraints, assumptions, acceptance conditions, and other task requirements into executable and verifiable commitments. It records the intended result, applicable constraints and their priorities, assumptions, conditions for proceeding, and other task information needed to construct a plan. User-supplied information and system inferences remain distinguishable, and unresolved assumptions stay attached to the commitments they affect. Each commitment can therefore be related to the subgoals and conditions that it governs. A change in a high-level commitment may alter several dependent subgoals, while a local change may affect only the connected part of the plan. New task information can update the specification and, in turn, alter dependent subgoals or conditions for proceeding.

NEWTON identifies a specification bottleneck in physics-grounded video generation: the physical specification must provide sufficient, dynamically relevant, and verifiable information for the simulator and generator to realize the intended behavior \cite{feng2026newton}. Divide and Conquer makes information including the objects, attributes, relations, and scene layout of a compositional request explicit before composing the image \cite{wang2024divide}. MetaPoint converts a free-form image-editing request into structured object-level commands, including precise spatial information for the requested edit \cite{zhou2026metapoint}.

\paragraph{Plan construction.}
Planning converts the current specification into an ordered and dependency-aware set of subgoals and intended actions. A plan states which subgoal is addressed, what must precede it, which information is required, and which condition must hold before a dependent subgoal can be addressed. It connects each subgoal to the action intended to address it and records the order in which related commitments are handled. The plan can cover the full trajectory or the remaining part of an ongoing task, depending on how much of the task has already been resolved. It may be represented in natural language, as a program-like structure, as a hierarchy or tree, or in another task-specific form, depending on the task and its available action interfaces. Across these forms, the relevant record is the organization of commitments into records that guide subsequent decisions; the representation determines how order, dependency, and alternative continuation are exposed.

GenArtist represents generation and editing operations in a planning tree, with sibling nodes encoding alternative tools for the same action \cite{wang2024genartist}. LightVA recursively decomposes an analytical request into data and visualization tasks for subsequent execution \cite{zhao2024lightva}. From Plans to Pixels generates structured atomic decompositions for open-ended, long-horizon image-editing requests and conditions tool and region selection on the resulting plan \cite{avg260515181}. Their plans differ in representation, while each records subgoals and the operations needed to address them.

\paragraph{Plan revision.}
Planning remains revisable when new user information or task evidence changes a commitment, a dependency, or another condition governing the expected next action. Such information may clarify an unresolved requirement, change a reference, introduce a new constraint, show that the current conditioning is insufficient for the intended result, or otherwise change the current specification. Revision compares the updated specification with the remaining plan, retains commitments that remain valid, and replaces the affected subgoals or their order. The scope of the revision follows the changed commitment: a local change can update the connected subgoals, while a change that governs several dependencies can require a broader reorganization of the remaining plan. This keeps the relationship between new information and the resulting planning decision explicit. Adaptive Task Reformulation changes the task representation and selects a new reformulation route when direct image editing fails \cite{avg260415917}. IMAGAgent orders constraint-dependent editing sub-tasks and uses reflection on intermediate results to trigger self-correction or retry \cite{avg260329602}. MIRA predicts one atomic editing action from the current visual state, feeds the updated image back into the next decision, and includes a stop decision when the instruction is satisfied \cite{avg251121087}.

Taken together, these operations form a progression from task interpretation to operational commitments, from commitments to dependency-aware subgoals, and from new information to an updated remaining plan. The resulting records provide the planning information required by later operations while keeping unresolved assumptions and plan changes traceable to the task information that produced them.



\subsubsection{Memory}

\paragraph{Memory definition.}
Memory is information retained from an interaction so that a later decision can use it with the current task information. A record becomes memory when the system selects it from an interaction, gives it a readable and addressable representation, and makes it available at a later decision point. The retained information can describe a requirement, reference, visual observation, artifact version, action outcome, or prior experience. Its representation can be textual, visual, structured, executable, or multimodal, with links that preserve the relation between a record and the operation or result from which it came. Interaction history and current context provide source information; memory is the selected record that remains available for later use. This distinction matters in visual creation because a later operation may need to recover the approved appearance of an entity, the source representation of an artifact, the evidence behind a previous decision, or the outcome of an earlier attempt.

Qwen-Image-Agent treats memory as a context-grounding source for historical and personalized information during image generation and editing \cite{avg260626907}. The Cognitive-structured Multimodal Agent externalizes visual information into episodic visual memory and uses a retrieval engine to reactivate relevant visual episodes across turns \cite{avg260708497}. LAVE maintains a memory buffer of prior conversations and combines recent history with the current instruction when constructing a video-editing action \cite{wang2024lave}. These systems show the same basic operation at different scopes: information produced earlier is selected and supplied to a later generation decision.

\paragraph{Memory formation, update, compression, deletion, and retrieval.}
Memory formation selects information with potential future utility from a request, reference, action, artifact output, observation, feedback signal, or user correction. The selection may preserve an event, extract a proposition from it, summarize several related events, or link an editable source to the result it produced. The selection criterion is the information's expected value for a later decision, together with its scope, reliability, and provenance. For example, during a multi-turn image edit, the system can retain a user-confirmed subject identity and a region that has already met the task constraints, while leaving an unconfirmed stylistic preference marked as unresolved. During long-video creation, frame-level observations can be collected as candidates and converted into a record about an entity or relation that persists across the relevant part of the trajectory.

Updating incorporates new information into an existing record or set of records. It can append a new event, revise a value, attach evidence, change the validity interval, merge compatible entries, or preserve competing versions when the evidence is unresolved. An update should keep the record's scope and source visible so that later decisions do not treat an obsolete observation as a current fact. If a later edit changes the approved appearance of an object, the corresponding memory can point to the new version while retaining the earlier version as provenance; if a new observation conflicts with an earlier one, the record can retain both claims with their conditions until a verification result resolves the conflict.

Compression reduces the amount of stored information or the context supplied to a decision while preserving the content needed for that decision. Summarization can replace repeated observations with a stable description, deduplication can merge records that express the same fact, and hierarchical storage can keep a compact summary linked to more detailed evidence. Compression is useful when a visual trajectory contains many frames, intermediate renders, or repeated tool results. A long sequence of consistent character observations, for instance, can be represented by a compact identity record linked to the frames that support it, while an exceptional frame remains available for diagnosis. The compressed record must preserve the constraints, dependencies, uncertainty, and provenance required to recover or revise the relevant part of the artifact.

Deletion removes a record, marks it inactive, or lets it expire when its validity, scope, utility, safety, or maintenance cost no longer supports retention. Deletion can be local: an obsolete tool parameter or a failed intermediate version may be removed while the task requirement and the last verified artifact remain available. It can also be conditional: a record may be retained for provenance but excluded from ordinary retrieval after its validity interval ends.

Retrieval makes retained information available to a current decision. The query is constructed from the task specification, the current observation, and the information needed for the next operation. Selection can use semantic match, entity or region identity, temporal scope, dependency links, recency, reliability, and modality; post-retrieval processing can order, filter, summarize, or bind the records to the current artifact. For example, when a later shot is edited, retrieval can return the latest approved identity record, the relevant reference frame, and the prior repair that affected the same dependency, while leaving unrelated experiments out of the decision context. Retrieval can occur at task initialization, at selected turns, or when a new observation creates a need for earlier evidence. Its role is complete when the returned record is available to the subsequent decision, such as choosing an input, editing a source representation, selecting a tool, revising a plan, or stopping.

COMFYCLAW distills trajectories, execution errors, and verifier feedback from image-workflow construction into reusable Agent Skills and progressively discloses those skills when later workflows require them \cite{avg260701709}. DataEvolver converts rejected text-rich image samples and their verification causes into semantic feedback, retains useful feedback entries, merges near-duplicate observations, and uses the resulting experience memory to revise later retrieval queries and generation prompts \cite{avg260631537}. VideoWeaver stores process and output evidence from long-video generation cases and uses the accumulated context to summarize category-level patterns and evolve reusable composition skills \cite{avg260608091}. These systems instantiate selected parts of formation, updating, compression, retention, and retrieval at different scopes.

\paragraph{Memory types.}
Memory types describe different properties of a retained record. The temporal dimension separates short-term memory, which supports decisions within a turn or current trajectory, from long-term memory, which remains available across turns, scenes, or tasks. The functional dimension distinguishes working memory for the active task workspace, factual memory for relatively stable information about users, environments, references, or visual entities, episodic memory for particular attempts and outcomes, and procedural or skill memory for reusable action patterns. These roles can be implemented by the same store when their formation and retrieval schedules differ.

The representation dimension concerns the carrier of the record. Textual or token-level summaries can preserve requirements and prior decisions; visual or multimodal records can preserve appearance, identity, and visual evidence; structured records and graphs can expose entities and dependencies; executable or artifact-linked records can connect source code, workflows, or other editable representations to their outputs. Access can be private to one role or shared among several roles, while the information channel can be textual, visual, structured, executable, or multimodal.

DreamFactory explicitly separates short-term information within a scene from long-term information that must persist across scene transitions, extracting a Base Description containing style, background, and character attributes and combining it with preceding-frame context for later key-frame generation \cite{xie2024dreamfactory}. ViMax uses global narrative context, retrieval-augmented generation, and graph-based cross-shot dependencies to carry story information into local video-generation decisions \cite{huang2026vimax}. ShareVerse implements spatiotemporal memory retrieval and conditions later video frames on retrieved memory frames to preserve a shared environment across asynchronous agent trajectories \cite{avg260302697}. Together, these examples cover temporal, functional, representational, and shared-access dimensions.

\paragraph{Functions in agentic visual generation.}
Memory in agentic visual generation serves three core functions: task continuity, episodic reuse, and procedural reuse. Task continuity preserves requirements, references, identities, dependencies, artifact versions, and decision evidence that remain relevant as a visual creation trajectory proceeds. A later operation can therefore continue from an established task interpretation, inspect an earlier source or artifact version, preserve completed work, and revise the unresolved part of the trajectory. In a multi-turn image edit, a retained subject record and the latest verified image can support a new local edit without losing the decisions established in earlier turns. In long-video creation, retained entity and scene records can carry an established relation into a later segment while local content changes.

Episodic reuse retrieves a record of a particular earlier attempt, outcome, failure, or repair for a new but related request. Such a record is useful when its task conditions, selected actions, resulting artifact, and review evidence remain interpretable. A system can retrieve a prior trajectory for a related visual defect, use the successful intervention as an initial candidate, and adapt it to the current artifact and constraints. It can also retrieve a failed episode to avoid repeating an action that produced the same conflict or to select an observation that distinguished the earlier failure from a successful case.

Procedural reuse stores a reusable action pattern distilled from one or more episodes. The pattern can specify an operation order, preconditions, parameter relations, expected observations, or a repair strategy while leaving task-specific values to the current request. Retrieved procedural information can guide tool selection, source-level editing, inspection order, or continuation decisions. For example, a reusable visual-document skill can preserve the sequence of inspecting source structure, rendering the result, checking the page, and revising the responsible source; a video-composition skill can preserve continuity checks while allowing the current scenes and assets to change.

Agent Banana uses Context Folding to abstract a growing editing history into asset-, execution-, and planning-level records; its persistent ActionContext keeps the verified effective editing path and the associated image states for later turns, providing task continuity across successive edits \cite{ye2026agentbanana}. Generation Navigator represents each subsequent action using the original prompt, prior actions, generated images, reviewer feedback, and scores accumulated along the trajectory, keeping task-local evidence available for the next decision \cite{avg260517969}. SEAR builds a self-evolving episodic memory that consolidates high-reward restoration trajectories into retrievable expertise and uses a state fingerprint to select relevant prior trajectories, supporting episodic and procedural reuse across related restoration requests \cite{avg260628971}.



\subsubsection{Tool}

This component represents a visual-creation operation as a callable tool and records the result of invoking it. The input to a tool call is constructed from the current task specification, the available task state, and the operation selected for the current subgoal. The tool boundary determines which inputs can be supplied, which artifact or environment is affected, and which outputs are returned to the ongoing trajectory. A returned result can contain a new or modified artifact, a structured state, an execution status, an error, or other information exposed by the executor. The interface may also expose preconditions, side effects, latency, cost, or reversibility, among other interface properties, when the implementation defines them.

\paragraph{Tool definition.}
A tool is a bounded, callable capability that carries out an operation on a visual artifact, an editable representation, or a task environment. Its definition connects an operation name or endpoint to an input representation, an execution procedure, and a return representation. The input can be a prompt, image, mask, structured record, program, parameter set, or another task-specific value. The output can be a rendered result, an edited source, a geometric object, a chart, a document, an interface state, a simulator trace, an execution message, or another task-specific result. For example, an image-editing tool may accept a source image, a mask, and an edit instruction, invoke an editor, and return an edited image with the version or file identifier needed for the next call. A CAD tool may accept a parametric program and return geometry together with compiler or kernel status. A useful interface keeps the relationship between the supplied input and the returned result addressable, so a later operation can identify the artifact version or environment state produced by the call. The interface can expose validation rules and failure states at the boundary, while the interpretation of whether a result satisfies the task remains a separate operation.

GenClaw provides an example of code-driven image operations in which a visual instruction is translated into executable operations over a canvas representation \cite{ye2026genclaw}. TOOLCAD formulates text-to-CAD interaction around tool calls to a CAD engine and uses trajectories containing tool execution outcomes \cite{gong2026toolcad}. CAD-Assistant connects a vision--language model to executable CAD functions so that language decisions produce CAD operations and returned model states \cite{mallis2024cadassistant}. These systems expose different interfaces, while each makes the operation boundary part of the available execution space.

\paragraph{Tool types.}
Tool types describe the operation supplied by an interface and the substrate in which it executes. For agentic visual generation, the working groups include visual generation and editing tools; visual perception and analysis tools; retrieval and external-information tools; code, application, and rendering execution tools; and simulation or task-environment tools, among related tool functions. The groups describe the function provided to the current trajectory. One interface can combine several functions, such as an application tool that accepts editable code, renders a visual artifact, and returns execution messages. A perception or analysis tool contributes an output such as a measurement, label, relation, data summary, or another task-relevant signal. That output can inform a later decision when it is related to the current task requirements and available state.

Visual generation and editing tools produce or modify the artifact used by the task; a call may map a prompt and reference image to a new image, or a source frame and edit region to a revised frame. Visual perception and analysis tools expose properties of an artifact or its source; examples include returning detected objects, text, measurements, or a data summary. Retrieval tools supply references, examples, or external information that can be used to construct later inputs. Code, application, and rendering tools execute programs or operate within an application environment; a plotting tool can turn a table and Python code into a figure, while a presentation tool can turn source code into rendered pages. Simulation tools execute a candidate representation under a modeled environment and return a trace or status that can be passed to a later decision; a physics tool can receive an event program and return simulated positions, contacts, or an execution failure. AMACE uses separate agents for chart-code generation, chart rendering, and chart-quality assessment \cite{namgoong2025amace}. PlotGen combines code generation with numeric, textual, and visual feedback agents for scientific visualization \cite{goswami2025plotgen}. METAL separates chart generation, visual critique, code critique, and revision in a chart-to-code workflow \cite{li2025metal}. These examples show how tool type is determined by the operation and execution substrate, while the same system can expose several related interfaces.

\paragraph{Tool creation.}
Tool creation constructs a callable capability that can be registered and invoked by a later decision. The construction process may generate an implementation, wrap an existing model or application programming interface, compose several primitive operations, or attach an input--output interface to an executable procedure, among other constructions. It receives the intended operation, its dependencies, and the values that must cross the interface. It produces a callable object together with the information required to invoke it, such as accepted arguments, return format, environment requirements, and reported execution states, together with related interface metadata. Tests, example calls, or other checks can establish that the new interface is executable before it is placed in the available action space. A generated image, a source program, a workflow description, or a memory record can contribute to this process when the system uses it to implement, parameterize, document, or test a callable operation.

Tool creation can occur before a task, when a system prepares a library of domain operations, or during a task, when an unavailable operation is assembled from accessible components. For example, a system can wrap a renderer as a callable function that accepts source code and asset paths, runs the renderer in a specified environment, and returns page images plus an execution status. A newly created tool can be specialized to the current artifact, generalized for later requests, or composed into a higher-level operation, among other scopes. Its utility depends on the compatibility between its input and output representations and the state of the execution environment. The record should therefore preserve the implementation version, dependencies, invocation conditions, and returned status needed to reproduce a later call. Tool creation supplies an additional callable capability; selecting when to use it and interpreting its result are separate decisions.

\paragraph{Tool selection and routing.}
Tool selection identifies the capability used for the current subgoal. The selected capability may be a generator, editor, analysis function, retrieval service, application, simulator, or another available executor. Selection uses the task requirements, current artifact and state, interface compatibility, and operating constraints such as quality, cost, latency, or side effects. Its output is a tool choice together with the arguments and conditions needed for invocation.

Routing describes how the task or intermediate result is sent through the selected capability or a sequence of capabilities. A route can connect one tool directly to the next operation, pass an artifact through several compatible tools, or direct different subgoals to specialized paths. Selection therefore answers which capability is used; routing describes the path, ordering, and handoff through which that capability participates in the task. Both can be established from the initial request and can be revised when an artifact, execution status, error, or other runtime information changes the requirements for the next operation. 

\paragraph{Tool use across artifact, execution, and coordination settings.}
Tool use is the invocation of a selected capability and the handoff of its returned result to the next operation. The handoff preserves the representation needed by the recipient: an image or video render can be passed as a visual artifact, source code as an editable representation, a chart table as structured data, a CAD program as executable input, a simulator trace as an environment result, or another task-specific representation. The return record can include success or failure status, changed files or objects, generated artifacts, execution metadata, or related execution information exposed by the interface. A later operation can then continue from the returned state, supply a revised input, select another compatible tool, or deliver the artifact to the user.

The first dimension is the artifact being produced or edited. Pixel and temporal artifacts include images, masks, reference frames, clips, and scene segments. Structured artifacts include source code, layer hierarchies, scene descriptions, geometric programs, data mappings, layouts, and interface states. A portrait-editing call may combine a source image, a subject mask, and an edit instruction, then return a new image and an editable region for another local change. A CAD call may pass a parametric program and return geometry with compiler or kernel status. A chart, presentation, or Web call may pass data and source code and return a figure, rendered pages, or an updated interface state.

The second dimension is the execution substrate. A selected tool can call a model or API, run code in a runtime or sandbox, operate a desktop or Web application, invoke a renderer, query an external information service, or execute a simulator or other task environment. A code-mediated operation can run a plotting script and return a figure plus execution status. An application-mediated operation can modify a slide or interface and return rendered pages or an updated state. A simulator-mediated operation can receive an event program and return a trace, contact states, or an execution failure. The handoff records the input representation, the executor, and the returned artifact or status so the next operation can address the correct file, object, or environment state.

The third dimension is the organization of tool calls. One decision process can construct inputs and invoke tools directly. A staged workflow can pass the output of one operation to a later operation with a different interface. A multi-agent workflow can assign generation, editing, programming, rendering, or integration calls to different roles. Mora organizes model and agent roles for text-to-video, image-to-video, and editing workflows \cite{yuan2024mora}. GenMAC distributes compositional text-to-video work among collaborating agents \cite{huang2024genmac}. In such workflows, a prompt interpretation can be handed to a generation role, a generated result can be handed to an editing or integration role, and a rendered output can be handed to a later delivery step. The execution record retains the input, the acting role and tool, and the returned artifact or status so that each handoff remains traceable.

These dimensions can be combined in the same trajectory. A task may begin with a model call over a reference image, move to code execution for a structured layout, pass the result to an application renderer, and then use a simulator or interaction environment for a task-specific check. Across these combinations, the tool-use record identifies the operation invoked, the representation crossing the interface, the executor or role that acted, and the result made available for the next decision.


\subsubsection{Perception}

This component obtains and interprets information about a visual artifact, its editable representation, its execution process, the task environment, references, and interaction. Its input is a task-relevant signal source; its processing acquires, extracts, compares, localizes, or checks information; and its output is structured evidence that a later Action can use. This evidence may contain a score, a pass or graded status, a failed requirement, a location, a likely cause, uncertainty, supporting observation, or a candidate target. The mechanism can be instantiated for different visual artifacts and execution settings.

\paragraph{Perception types and roles.}
Perception can operate on several kinds of information. Appearance perception concerns visible content, such as object presence, text rendering, color, texture, and layout. Structural perception concerns the representation that produced the visible result, such as source code, layers, scene relations, geometry, data mappings, or interface structure. Temporal and relational perception concerns changes across frames, shots, turns, entities, or dependent elements. Execution and environment perception concerns compiler messages, runtime status, application state, browser state, simulator traces, or other signals returned by an external process. Reference and user perception concerns comparisons with a supplied image, specification, example, preference, or correction. For example,in a slide workflow, the trajectory may combine page appearance, source structure, and parser status; a CAD workflow may combine rendered geometry, kernel measurements, and program state; a video workflow may combine sampled frames with identity and temporal relations.

\paragraph{Signal acquisition and observation.}
Signal acquisition obtains information from the artifact and its execution context. A system may, for instance, sample outputs at a fixed cadence, check a known risk point, select a crop or frame range, inspect source code or geometry, query a data record, retrieve a reference, or read a compiler, renderer, browser, application, or simulator result. The useful properties of a channel include its coverage, timing, localization, and cost. For instance, a scheduled render check may follow a scene transition, while a local inspection may be selected when uncertainty is concentrated in one region. These properties shape which signals are collected and when they become available for interpretation.

In a program-mediated video workflow, Kubrick inspects Blender renders within its synthetic-video process \cite{he2024kubrick}. The artifact and execution substrate determine which signals become available for interpretation.

\paragraph{Interpretation and diagnosis.}
Interpretation connects an acquired signal to a task requirement, a possible violation, and an affected part of the trajectory. Diagnosis can identify a region, frame, object, component, or operation; associate the observation with a likely source; and retain competing hypotheses or confidence when the system reports them. The usefulness of a diagnosis depends on how directly it connects the observed signal to a later action target. 

Concrete cases include CADSmith, which combines execution checks, CAD-kernel measurements, and visual assessment to separate runtime, dimensional, and morphological problems \cite{barkley2026cadsmith}; IterCAD, which uses executable sandbox feedback and geometric evidence in interactive CAD generation and editing \cite{hu2026itercad}; PhysAgent, which uses stage-specific simulation results to distinguish reconstruction failures from dynamics failures \cite{li2026physagent}; and DirectorBench, which organizes long-form video evaluation around checkpoint-level specifications and profile-aware bottleneck reports \cite{avg260530090}. These cases treat diagnosis as an interpretation of available evidence; the repair, retry, or route change remains a subsequent Action.

\paragraph{Verification and acceptance evidence.}
Verification checks an explicit requirement, acceptance condition, or property of the current result. Relevant targets include perceptual and semantic correspondence; relational, spatial, and geometric validity; temporal continuity and identity consistency; data, factual, and numerical fidelity; executable behavior and structural validity; physical plausibility; and human acceptance when a human review is part of the reported system. For a text-rendering requirement, this may involve OCR; for a CAD constraint, a geometry or compiler check; for a physical event, a simulator; and for an open-ended visual property, a learned judge or human review. Different targets require different evidence, and each check exposes a particular uncertainty.

A verification process can combine exact tests, reference-based comparisons, learned judges, source-level checks, simulations, and human review. It can record coverage, granularity, latency, confidence, and disagreement, and can schedule cheap checks before expensive tests.  These checks answer different questions about the current result and can be combined when the artifact exposes multiple representations.


\subsubsection{Action}

\paragraph{Action definition.}
Action is an operation selected at a decision point that changes a visual artifact, editable representation, external environment, available information, interaction history, or control trajectory. It receives the current goal and specification, Memory records, artifact and environment state, available Tool interfaces, and Perception evidence. It constructs an input, invokes an operation, and incorporates the returned artifact, state, status, communication result, or control update into the running trajectory. An action record can therefore identify the selected operation, its input and target, its executor or role when reported, its returned result, and available cost, latency, reversibility, or side-effect information. The action space includes artifact transformation, information acquisition, execution, environment interaction, communication, coordination, and trajectory control, among related operations.

The same action principle applies across several operating arrangements. A single decision process may invoke a model or editor directly; a staged workflow may pass a source representation through generation, rendering, and inspection; and multiple roles may divide generation, editing, programming, integration, or delivery. For instance, a reference image can be passed to a generator, the generated representation to a code or application runtime, and the rendered result to a later operation.

\paragraph{Artifact actions.}
Artifact actions create or modify a visual artifact or its editable representation. The mechanism maps a task-specific input to a changed artifact while preserving the representation needed for a later operation. It can modify pixels, frames, masks, source code, layers, geometric programs, data mappings, layouts, or executable interfaces. A prompt and reference image can produce an image; a source frame, region, and edit instruction can produce a revised frame; a parametric CAD program can produce geometry and execution status; and data with source code can produce a figure, transformed data, or rendered pages. The operation may be global or localized, and its granularity determines which completed properties remain available for later steps.

T2I-Copilot combines structured prompt interpretation, model routing, segmentation, interactive canvas operations, and quality evaluation in a text-to-image workflow \cite{chen2025t2icopilot}. AutoStudio separates prompt extraction, subject management, layout construction, layout supervision, and drawing for multi-turn image generation \cite{cheng2024autostudio}. Anywhere assigns distinct operations to foreground understanding, diversity enhancement, object-integrity protection, and text consistency in foreground-conditioned image generation \cite{xie2024anywhere}. The examples differ in artifact and operation, while each makes an artifact change available to later steps.

\paragraph{Information acquisition actions.}
Information acquisition actions obtain material for a later operation. The mechanism constructs a query or inspection request, sends it to a source, and returns a reference, measurement, source fragment, preview, execution message, or other task-relevant record. It includes search and retrieval of references or external knowledge, data queries, source-code and file inspection, browser or application inspection, environment-state inspection, rendering of crops or previews, selection of a frame range, invocation of a perception or analysis service, and read or write operations on retained task records. For instance, a system may retrieve a visual reference before generation, inspect a source file before editing, or render a local crop before selecting a region-level action. Perception interprets the returned material; Action performs the operation that obtains it.

\paragraph{Execution and environment actions.}
Execution and environment actions operate through a model or API, a code runtime or sandbox, an application or browser, a renderer or compiler, a simulator, or another task environment. The mechanism translates a selected operation into an executable call, applies it to a target artifact or environment, and returns the resulting artifact, state, status, error, or trace. A CAD action can submit a program to an executable engine, an application action can modify a source file and render a page, and a simulator action can execute an event program and return a trace or failure. The Tool section defines the callable boundary; Action concerns the operation performed through that boundary and the state change it produces.

\paragraph{Communication and coordination actions.}
Communication and coordination actions move information between roles or authority boundaries. A sending operation selects what information to transmit, formats it for the recipient, and records the response or resulting handoff. The action can involve delegation, message passing, synchronization, aggregation, role handoff, delivery of an artifact or structured state, clarification requests, approval requests, rejection handling, progress reporting, or human override. In a multi-agent workflow, a prompt interpretation can be sent to a generation role, a generated result can be sent to an editing or integration role, and a rendered output can be sent to a later delivery step. Mora organizes model and agent roles for text-to-video, image-to-video, and editing workflows; GenMAC distributes compositional text-to-video work among collaborating agents \cite{yuan2024mora,huang2024genmac}.

\paragraph{Evidence-conditioned action and trajectory control.}
Evidence-conditioned action uses Perception output to determine how the trajectory continues. The mechanism receives a score, failed requirement, location, likely cause, or uncertainty, maps it to an admissible response, and performs that response. A critique can be sent to a prompt or code generator, a failed requirement can be applied to a local edit or revised input, an execution error can be passed to a retry or routing decision, and a result can be communicated to another role or user. A verified state can support continuation; an unresolved state can lead to more evidence, clarification, recovery, or stopping.

Trajectory-control actions include tool, model, or role selection and routing; prompt, input, plan, dependency, or parameter updates; retry, branching, continuation, or path substitution; local repair, replanning, rollback, and bounded relaxation; stopping, termination, or deferral; and human authorization, override, or escalation. CoSTA* formulates multi-turn image editing as cost-sensitive tool-path selection \cite{gupta2025costa}. FaSTA* separates fast and slow tool paths and mines reusable subroutines for later editing decisions \cite{gupta2025fasta}. 


\subsubsection{Cross-Task Self-Improvement}

\paragraph{Definition and scope.}
Cross-Task Self-Improvement is a persistent change in a system's later-task decision process that is produced from experience collected in completed trajectories. Persistence means that the retained information, learned procedure, routing behavior, verifier behavior, or model state survives the completed trajectory and remains available when a subsequent task begins. The later decision can concern reference retrieval, prompt construction, tool selection, action ordering, verification, repair, or another task operation. Action covers corrections performed during the active trajectory. Memory covers information retained for later access. Cross-Task Self-Improvement covers the update that makes retained experience alter a later task decision. Within-action correction operates inside one selected operation. Within-trajectory iteration uses observations and feedback before the current request ends. Episodic cross-task reuse retrieves a prior case or trajectory for a subsequent request. Persistent system adaptation changes a reusable skill, routing rule, construction policy, verifier, or model behavior across subsequent requests. Final artifact quality is an outcome measure. Candidate selection is a decision inside the active trajectory. A later-task behavior change is the defining unit of cross-task self-improvement.

\paragraph{Experience formation and attribution.}
Experience formation converts a completed trajectory into information that can guide later decisions. The process records the task specification, references, intermediate artifacts, editable representations, tool calls, execution status, observations, diagnoses, verification results, critiques, repairs, costs, and user or evaluator judgments, together with provenance and task context. The formation process selects records with future decision value and compresses them into a case, a failure explanation, a repair relation, a preference, or another usable learning signal. Attribution links the observed outcome to the actions, conditions, and evidence that preceded it. This link attributes the observed improvement or failure to earlier choices, tool paths, repairs, and their conditions, and identifies the part of that relation that can be reused on a later task. The mechanism can compare successful and unsuccessful trajectories, align intermediate evidence with later outcomes, and discard records whose quality or provenance is unresolved.

GenEvolve compares multiple tool-orchestrated image-generation trajectories for the same request and extracts best--worst differences into structured visual experience covering search, reference selection, prompt--reference construction, and failure avoidance; the experience is supplied to a teacher branch that distills guidance into a student policy \cite{chen2026genevolve}. SIDiffAgent records complete diffusion-agent trajectories, summarizes pitfalls and successes at decision nodes, retrieves similar records, and injects corrective and workflow guidance into later prompts; its second episode uses the accumulated database on later benchmark prompts \cite{avg260202051}. These systems show experience formation as a sequence of logging, outcome attribution, abstraction, and later guidance construction.

\paragraph{Reusable representations.}
Reusable experience has several representation forms with different operational roles. A case preserves an episode, an artifact, a repair, or a tool trace for retrieval or replay on a related request. A strategy compresses recurring relations into a rule about search, reference choice, prompt construction, verification, or repair order. A skill packages a callable procedure, code fragment, workflow, or interface with inputs and execution conditions. The representation controls the amount of context carried into a later task and the degree of direct execution available to the system. Case reuse retains task-specific detail. Strategy reuse retains an abstraction that can guide several related tasks. Skill reuse exposes an executable operation that can be invoked as part of a later workflow.

A long-video composition procedure can specify how clips, references, timing, and continuity checks are assembled. An image-editing correction rule can map a recurring defect to an edit operation and its parameters. A data-construction feedback rule can convert repeated rejection causes into the retrieval and synthesis choices for the next construction round. VideoWeaver refines and merges composition and creator skills from execution-trace and final-video evaluation, then tests the evolved skills on held-out task categories \cite{avg260608091}. COMFYCLAW distills trajectories, execution errors, and verifier feedback into a progressively disclosed skill library for ComfyUI workflow construction and validates the library on held-out prompts \cite{avg260701709}. The two systems instantiate skill-level reuse with different visual products and workflow interfaces.

\paragraph{Retention, retrieval, and adaptation locus.}
Retention representation, access mechanism, and adaptation locus describe different parts of the update process. Retention representation specifies what survives a trajectory, such as an episode, summary, rule, skill, workflow, feedback record, or parameter state. Retrieval specifies how a later task queries that information, using signals such as task similarity, degradation-aware state fingerprints, artifact properties, tool compatibility, or temporal scope. Adaptation locus specifies what changes after retrieval or learning: the prompt, tool route, plan, verification rule, construction policy, skill library, or model parameters. An episode can therefore be stored in memory, retrieved by a state fingerprint, and applied as a prompt-construction change. A skill can be stored in a library, invoked through an interface, and used to alter tool ordering. A routing record can be retrieved from prior tool outcomes and used to select a generator.

OctoT2I updates a text-to-image router through repeated evaluation and learning, retaining tool-use knowledge that changes later routing and supports new-tool integration \cite{jiang2026octot2i}. EvoIR-Agent organizes restoration experience in a hierarchical pool and retrieves it at multiple levels to guide later diagnosis and tool ordering \cite{avg260522208}. SEAR indexes explored restoration trajectories with degradation-aware state fingerprints and reuses the resulting episodic memory to reduce later search cost \cite{avg260628971}. DataEvolver retains verifier scores and rejection causes, summarizes them as semantic feedback, and feeds the updated feedback memory into later retrieval and targeted synthesis; the resulting construction data transfers to a second evaluated generator \cite{avg260631537}. These cases place the update at different loci while preserving the same sequence of retention, access, and later decision change.

\paragraph{Transfer and evaluation.}
Evaluation determines whether a retained experience changes a later task and whether the change remains useful under new conditions. In-task improvement measures the quality change across revisions of one request. Cross-task transfer measures later requests after experience construction. Generalization tests unseen prompts, task families, tool combinations, generators, executors, or related environments. Persistence tests whether the update remains available after the original trajectory and across later sessions or batches. Safety evaluation measures forgetting, harmful updates, evaluator overfitting, retrieval contamination, and regressions on previously verified capabilities.

A causal comparison holds fixed the factors outside the claimed adaptation mechanism and varies the retained experience or update. The trace records whether the later run retrieves a different reference, selects a different tool or order, constructs a different prompt, changes a verification threshold, chooses another repair, or takes another related decision. Artifact evaluation measures constraint satisfaction, semantic alignment, structural validity, visual quality, and domain-specific properties. Process evaluation records calls, latency, cost, retries, recovery steps, and stopping decisions. Held-out task families, unseen tool combinations, independent evaluators, and human judgments test transfer beyond the trajectories used to build the update. Ablations with matched generators and sampling budgets separate retrieval effects, extra sampling, generator strength, and persistent policy or skill changes.

VISTA supplies a test-time boundary case in video generation. Its procedure builds a temporal plan, generates candidate videos, selects a candidate through pairwise comparison, collects visual, audio, and contextual critiques, and rewrites the prompt for later rounds; its evaluation compares successive rounds on single- and multi-scene tasks \cite{long2026vista}. The reported adaptation occurs inside the test-time process. Persistent cross-task evidence uses a later-task retention test, transfer split, forgetting check, provenance record, and reversible update procedure.
